\label{sec:conclusoes}

The message-passing encoders on the graph of the active channels, with the event-level scalars, reconstruct the vertex of \qty{10}{\MeV} deposits (a mean of \qty{530.1}{\pe} per test event) with a mean 3-D error of \num{25.95} to \qty{26.27}{\cm} and a combined RMSE of \num{18.55} to \qty{18.80}{\cm} for the four best. At matched capacity, the mean 3-D error of the \num{11} encoders varies by up to \qty{2.89}{\cm}, and these four cannot be distinguished on the \num{937} test events. GATv2-cascade, chosen for the lowest mean 3-D test error, has a per-axis RMSE, in the sensor frame, of \num{20.45}, \num{14.35} and \qty{20.46}{\cm} in $x$, $y$ and $z$.

On the same events and on clean light, GATv2-cascade has a mean 3-D error \qty{1.71}{\cm} lower than that of XGBoost and a combined RMSE \qty{2.30}{\cm} lower; it is better on the dim events and in the tail, but worse in the brightest quartile. In $z$, the calibrated barycentre has an RMSE \num{2.07} times that of GATv2-cascade because of the vertices near the end cap, where the error of the barycentre is systematic, a bias that APEX did not find in its simulation. Below $z = {}$\qty{20}{\m}, the barycentre stays close to GATv2-cascade (\num{23.29} and \qty{22.54}{\cm}), and its error is mainly statistical: with every photon detected, it falls to \qty{8.33}{\cm}. The combined RMSE of GATv2-cascade is \qty{10.6}{\percent} lower than that of the network of the previous stage, a value computed from the published per-axis values, which come from a different partition and from the geometry prior to the corrections; the difference is in $x$ and $y$, since in $z$ the value of that network, \qty{20.87}{\cm}, lies within the \qty{95}{\percent} interval of GATv2-cascade. At equal detected light with the average yield, this combined RMSE is \qty{51.0}{\percent} lower than the combined width of APEX (\qty{41.3}{\percent} with the minimum and \qty{28.0}{\percent} at equal energy), a root mean square of widths with no published definition, read from a plot of a different simulation. In~$z$, against \qty{20.46}{\cm} for GATv2-cascade, APEX has \qty{24.8}{\cm} with the average yield, above the \qty{95}{\percent} interval, and \qty{19.4}{\cm} with the minimum, within it; only at equal energy, where it detects \num{3.4} times as much light, is it better (\qty{13.7}{\cm}).

{\clubpenalty=10000 The error of GATv2-cascade concentrates in the dim events, far from the long walls (a~mean 3-D error of \qty{42.79}{\cm} in the dimmest quartile, against \qty{11.67}{\cm} in the brightest), but the effect of photon statistics was not separated from that of position. The results hold for one energy, one efficiency, one geometry, one photon draw per event and vertices on the training grid. Each encoder was trained once, without the hyperparameter search carried out for the boosters, with different batch sizes and numbers of steps; there was no ablation of the \num{33} scalars; the boosters were trained on augmented light; and the test events had already guided the final design of the encoders.\par}

{\clubpenalty=10000 As future work, this report proposes training and evaluating the encoders for several efficiencies, energies and geometries; generating several photon draws per event and vertices off the grid; repeating the comparison with a fixed optimisation budget, a hyperparameter search per encoder, boosters tuned and trained on clean light, and several seeds, with and without the scalars, on new test events; predicting the coordinates from the channel states, so as to test whether summarising the graph in a single vector limits the precision; combining boosters and encoders, which err less in different light ranges; comparing with APEX in an equivalent simulation, confirming the correspondence of the axes and the definition of the width; and measuring the inference time per event.\par}
